\pdfbookmark[0]{Beobachtungen der magnetischen Variation. 1834}{variation_1834_october}
\begin{center}\itshape
Beobachtungen der magnetischen Variation in Göttingen und Leipzig\\
am 1. und 2. Oktober 1834.
\end{center}
\begin{center}
\hrule\vspace{2mm}
\textsc{Poggendorff}. Annalen der Physik und Chemie. 1834. Bd.\ 33.
\vspace{2mm}\hrule\vspace{5mm}
\end{center}
Die in meinem Aufsatz über das hiesige magnetische Observatorium erwähnten Beobachtungen der magnetischen Variation an den verabredeten Tagen wurden seitdem hier noch zweimal angestellt, am 6. und 7. August und am 23. und 24. September. Beim ersten Termin traten recht starke und bemerkenswerte Anomalien auf. Umso mehr ist zu bedauern, dass zufällige Umstände entsprechende Beobachtungen an anderen Orten verhinderten. Die Septemberbeobachtungen wurden hingegen auch
\SourceLeaf{532}in Leipzig und Berlin ganz vollständig und in Braunschweig beinahe vollständig angestellt. Außerdem wurden sie zur Hälfte in Kopenhagen angestellt, wo man aus Versehen den 24. und 25. September statt des 23. und 24. gewählt hatte. Eine vollständige Veröffentlichung dieser Beobachtungen wäre allerdings weniger interessant, da der Verlauf an diesen beiden Tagen sehr regelmäßig war, obwohl mehrere an sich sehr kleine Anomalien in den ersten 24 Stunden an allen vier Orten eine bewundernswerte Übereinstimmung zeigten. Bemerkenswert bleibt jedoch, dass laut einer Zeitungsmeldung am Abend des 23. September in Glasgow ein sehr starkes Nordlicht gesehen wurde. Es hatte also ganz entschieden zumindest keinen Einfluss auf die Magnetnadel, der bis nach Norddeutschland reichte.

Die Anwesenheit von Professor \textsc{Weber} in Leipzig gab inzwischen Anlass, noch einige zusätzliche Stunden für gleichzeitige Beobachtungen in Göttingen und Leipzig festzulegen. Gewählt wurden am 1. und 2. Oktober morgens $7\tfrac12$ bis $8\tfrac12$, mittags $12\tfrac12$ bis $1\tfrac12$ und abends 8 bis 10 Uhr. Abgesehen von einigen kleinen Versäumnissen wurden diese Stunden an beiden Orten eingehalten. Im hiesigen magnetischen Observatorium beobachtete mein Sohn \textsc{Wilhelm Gauss}, in Leipzig beobachteten Professor \textsc{Weber}, Professor \textsc{Möbius} und Dr.\ \textsc{Thieme}. — — — Man wird mit Vergnügen die große Übereinstimmung sowohl bei den großen Bewegungen am Abend des 1. Oktober als auch bei fast allen kleinen Bewegungen bemerken. Daran ist unzweifelhaft zu erkennen, dass sie von Kräften ausgehen, die über große Entfernungen wirken, derzeit aber noch sehr rätselhaft sind. In Leipzig waren die Anomalien im Allgemeinen etwas kleiner als in Göttingen. Der Herd der wirkenden Kräfte wird also näher bei Göttingen gelegen haben. Ich bemerke nur noch, dass ich während eines Teils dieser Stunden selbst mit einem zweiten, in der hiesigen Sternwarte aufgestellten Apparat beobachtet habe, über den ich bald ausführlicher zu berichten gedenke. Diese Beobachtungen ergaben bei den größeren und kleineren Bewegungen einen fast vollkommenen Gleichlauf mit denen des hiesigen magnetischen Observatoriums. Ein ähnliches Ergebnis zeigte sich auch am 23. und 24. September und bei vielen anderen Versuchen. So konnten die Uhren an beiden Orten schon öfter allein anhand der magnetischen Erscheinungen bis auf einen kleinen Bruchteil einer Zeitminute genau verglichen werden. Dasselbe gelingt mithilfe
\SourceLeaf{533}der größeren Bewegungen am 1. und 2. Oktober zwischen Göttingen und Leipzig, wo die Uhren an beiden Orten nur wenige Sekunden von der mittleren Ortszeit abwichen.

Durch diese Erfahrungen gewinnen nun auch die kleinen, in sehr kurzen Abständen wechselnden Schwankungen der Magnetnadel außerordentliches Interesse. Es ist zu wünschen, dass auch sie durch Beobachtungen an vielen weit voneinander entfernten Orten sorgfältig verfolgt werden. Deshalb ist es unbedingt nötig, alle Beobachtungen in recht kurzen Zeitabständen vorzunehmen. Bisher beobachteten wir von 5 zu 5 Minuten. Aber auch dieses Intervall ist fast noch zu lang, und wir beabsichtigen, an den verabredeten Tagen künftig immer von 3 zu 3 Minuten den Stand der Magnetnadel zu bestimmen. Ich muss dabei erwähnen, dass wir das vom Herausgeber dieser Annalen (Bd.\ XXXII, S.\ 569 bis 572) erläuterte Verfahren nur anfangs gebraucht und schon lange durch ein etwas abgewandeltes ersetzt haben. Um den Stand der Magnetnadel für einen Zeitpunkt zu erhalten, beobachten wir sie zu sechs verschiedenen Zeitpunkten, die jeweils um die Dauer einer Schwingung auseinanderliegen, und zwar so, dass der gewünschte Zeitpunkt in die Mitte fällt. Anstelle der genauen Schwingungsdauer wählen wir die nächstliegende ganze Zahl von Sekunden, genauer von Uhrschlägen, zum Beispiel im magnetischen Observatorium $20''$ statt $20''4$. Die Beobachtungen am 2. Oktober für $8^{\mathrm h}\,15'$ abends lauteten daher wie folgt.
\begin{center}
\begin{tabular}{rrr|r}
$8^{\mathrm h}$ & $14'$ & $10''$ & $672.6$\\
 & & $30$ & $672.3$\\
 & & $50$ & $671.3$\\
 & $15$ & $10$ & $671.8$\\
 & & $30$ & $669.9$\\
 & & $50$ & $670.8$
\end{tabular}
\end{center}
\Needspace{13\baselineskip}
Daraus ergeben sich fünf Mittelwerte, die tatsächlich den daneben angegebenen Zeiten entsprechen.
\begin{center}
\begin{tabular}{rrr|r}
$8^{\mathrm h}$ & $14'$ & $20''$ & $672.45$\\
 & & $40$ & $671.80$\\
 & $15$ & $0$ & $671.55$\\
 & & $20$ & $670.85$\\
 & & $40$ & $670.35$
\end{tabular}\\[2mm]
\begin{tabular}{lrr}
und daraus der Mittelwert für & $8^{\mathrm h}\,15'$ & $671.40$
\end{tabular}
\end{center}

\SourceLeaf{534}Ich habe absichtlich dieses Beispiel gewählt, bei dem die Nadel schnelle Änderungen zeigte. Die selbst von 20 zu 20 Sekunden so deutlich hervortreten. Wir haben genügend Fälle, in denen sich ein ähnliches Ergebnis selbst in halb so großen Zeitabständen zeigt. Die Erklärung bestimmter Abwandlungen dieser Beobachtungsart, die wir öfter anwenden, und die Rechtfertigung dieser mit Bedacht gewählten Art der Mittelbildung muss ich mir für eine andere Gelegenheit vorbehalten. Ich muss aber erwähnen, da der Herausgeber dieser Annalen dies an der angeführten Stelle nicht ausdrücklich bemerkt hat, dass es eine wesentliche Bedingung für die Zulässigkeit all dieser Beobachtungsarten ist, die Nadel zuvor so weit wie möglich zur Ruhe gebracht zu haben, sodass die Schwingungen nur wenige Skalenteile betragen. Im hiesigen magnetischen Observatorium ist es im Laufe der Beobachtung selten nötig, die Nadel auf diese Weise zu beruhigen oder zumindest die Beruhigung zu wiederholen. Wer aber in einem weniger günstigen Raum beobachtet, darf keinesfalls versäumen, dies in der Zwischenzeit mit den bekannten Mitteln zu tun, sooft es nötig wird.

Da die Beobachtungen bei der jetzt als notwendig erkennbaren Verkürzung der Zwischenzeiten sehr viel mühsamer werden als früher, als nur von Stunde zu Stunde aufgezeichnet werden musste, wurde von mehreren Seiten der Wunsch geäußert, künftig sowohl die Anzahl als auch die Dauer der Termine etwas zu verringern. — — —

Göttingen, den 5. November 1834.

